% ================================================================
% Capitolul 6 – Evaluare
% Lucrare de licență: Kernel Trap – Sistem Inteligent de Prevenire
% a Intruziunilor cu Componente de Honeypot și Învățare Automată
% Autor: Banciu George-Horațiu
% Coordonator: Conf. Univ. Dr. Bufnea Darius Vasile
% ================================================================

\chapter{Evaluare}

Capitolul de față prezintă evaluarea experimentală a sistemului KernelTrap.
Sunt enunțate, mai întâi, asumările de securitate pe care se sprijină
funcționarea corectă a sistemului; sunt apoi raportate metricile de
performanță măsurate pe mediul de testare, urmate de trei scenarii concrete
de atac executate cu unelte folosite curent în comunitatea de securitate
ofensivă; secțiunea finală tratează limitările sistemului și clasele de
atacatori pe care îi acoperă cu eficacitate maximă.

% ----------------------------------------------------------------
\section{Asumări de securitate}
% ----------------------------------------------------------------

KernelTrap este un sistem de prevenire a intruziunilor proiectat pentru un
scenariu specific: un cont SSH neprivilegiat compromis pe o gazdă altfel
necompromisă. Validitatea protecției oferite depinde de două asumări fundamentale
asupra mediului de rulare, sintetizate în tabelul \ref{tab:cap6-trust}.

\begin{table}[h]
  \centering
  \caption{Asumările de securitate ale sistemului}
  \label{tab:cap6-trust}
  \renewcommand{\arraystretch}{1.08}
  \begin{tabular}{|p{0.32\textwidth}|p{0.58\textwidth}|}
    \hline
    \textbf{Asumare} & \textbf{Consecință} \\
    \hline\hline
    Contul \texttt{root} nu este compromis.
      & Agentul, scriptul de pivotare, regulile \texttt{iptables} și controlul containerului honeypot rulează cu privilegii ridicate. Un atacator care a obținut deja drepturi \texttt{root} poate dezactiva agentul, șterge regulile de firewall, opri procesul de scoring sau modifica însăși logica de pivotare. \\
    \hline
    Atacatorul nu are acces fizic la mașină.
      & Atacuri prin consola fizică (\texttt{tty*}), dispozitive USB ostile, modificare hardware, reset GRUB sau extragere de disc cad în afara domeniului de aplicabilitate. Filtrarea agentului pe sesiuni SSH \texttt{pts/*} reflectă explicit această alegere. \\
    \hline
  \end{tabular}
\end{table}

Din aceste asumări derivă două consecințe operaționale: telemetria colectată
prin eBPF/Tracee este considerată de încredere atâta timp cât kernelul
însuși nu este compromis; iar scriptul de pivotare, deși expus pentru a
putea fi rulat fără parolă de către agent (printr-o regulă \texttt{sudoers}
de tip \texttt{NOPASSWD} aplicată strict pe comanda \texttt{hp\_pivot\_user}),
nu este proiectat să reziste unui atacator care a obținut deja drepturi
de modificare a fișierelor de sistem.

% ----------------------------------------------------------------
\section{Mediul de testare și metodologie}
% ----------------------------------------------------------------

Configurația folosită pentru evaluare este descrisă în tabelul
\ref{tab:cap6-env}. Topologia a constat dintr-o gazdă monitorizată, serverul
central de analiză și containerul honeypot, rulate pe aceeași subrețea
locală pentru a izola măsurătorile de variabilitatea WAN.

\begin{table}[h]
  \centering
  \caption{Configurația mediului de testare}
  \label{tab:cap6-env}
  \renewcommand{\arraystretch}{1.08}
  \begin{tabular}{|l|p{0.55\textwidth}|}
    \hline
    \textbf{Componentă} & \textbf{Configurație} \\
    \hline\hline
    Gazda monitorizată       & TODO: <distribuție Linux, kernel, CPU, RAM>  \\
    \hline
    Serverul central         & TODO: <distribuție, RAM, CPU>                \\
    \hline
    Container honeypot       & Docker (Ubuntu 22.04), capabilități \texttt{AUDIT\_WRITE}, \texttt{AUDIT\_CONTROL} \\
    \hline
    Versiune Tracee          & \texttt{aquasec/tracee:latest} (TODO: <commit/tag>) \\
    \hline
    Versiune Python          & TODO: <3.x>                                   \\
    \hline
    Versiune Redis           & 7 (imagine oficială Docker)                   \\
    \hline
  \end{tabular}
\end{table}

Evaluarea a vizat două categorii de proprietăți: (1) overhead-ul de
performanță introdus pe gazda monitorizată, în repaus și în regim de lucru;
(2) capacitatea de detecție pe trei scenarii cu unelte publice de atac.

% ----------------------------------------------------------------
\section{Performanță}
% ----------------------------------------------------------------

Tabelul \ref{tab:cap6-perf} sintetizează metricile măsurate pe mediul descris
mai sus, raportate la țintele nefuncționale enunțate în capitolul 5
(tabelul \ref{tab:cap5-nfr}).

\begin{table}[h]
  \centering
  \caption{Metrici de performanță}
  \label{tab:cap6-perf}
  \renewcommand{\arraystretch}{1.08}
  \begin{tabular}{|l|l|l|l|}
    \hline
    \textbf{Metrică} & \textbf{Valoare} & \textbf{Țintă} & \textbf{Verdict} \\
    \hline\hline
    CPU agent (idle)                          & TODO: \texttt{<X\%>}  & < 2,5\% (NF1) & TODO \\
    \hline
    CPU agent                                 & TODO: \texttt{<X\%>}  & < 2,5\% (NF1) & TODO \\
    \hline
    RAM agent                                 & TODO: \texttt{<X> MB} & ---           & ---  \\
    \hline
    RAM server central                        & TODO: \texttt{<X> MB} & ---           & ---  \\
    \hline
    Throughput susținut server                & TODO: \texttt{<X> ev/s} & ---         & ---  \\
    \hline
  \end{tabular}
\end{table}

% ----------------------------------------------------------------
\section{Scenarii de atac și răspunsul sistemului}
% ----------------------------------------------------------------

Pentru a verifica capacitatea de detecție pe acțiuni cu profil realist, au
fost executate trei scenarii pe gazda monitorizată, după autentificarea unui
cont SSH neprivilegiat. Fiecare scenariu este descris prin succesiunea de
acțiuni ale atacatorului, indicatorii observați de KernelTrap și verdictul
sistemului.

\subsection{Scenariul 1: recon automat cu LinPEAS}

LinPEAS este un script Bash open-source folosit pe scară largă pentru
enumerarea automată a căilor de privilege escalation pe sisteme Linux. La
execuție lansează în cascadă zeci de procese pentru colectarea de informații
despre sistem, conturi, configurări de servicii și permisiuni.

\begin{table}[h]
  \centering
  \caption{Scenariul 1: indicatori observați la rularea LinPEAS}
  \label{tab:cap6-scen1}
  \renewcommand{\arraystretch}{1.08}
  \begin{tabular}{|p{0.42\textwidth}|p{0.48\textwidth}|}
    \hline
    \textbf{Acțiunea atacatorului} & \textbf{Indicator KernelTrap} \\
    \hline\hline
    \texttt{wget}/\texttt{curl} script de pe host extern   & \texttt{is\_network\_tool=1}, \texttt{has\_url\_or\_ip=1} \\
    \hline
    \texttt{id}, \texttt{whoami}, \texttt{uname}, \texttt{hostname}, \texttt{ss}, \texttt{ps} în cascadă & \texttt{is\_basic\_recon=1} repetat \\
    \hline
    Citire \texttt{/etc/passwd}, \texttt{/etc/shadow}, \texttt{.ssh/}      & \texttt{has\_sensitive\_path=1} \\
    \hline
    Rafale dense de syscall-uri          & \texttt{syscall\_rate\_5s} și \texttt{unique\_syscalls\_5s} ridicate \\
    \hline
  \end{tabular}
\end{table}

Rezultat: severitatea 2 apare pe majoritatea evenimentelor în primele
secunde de execuție; fereastra glisantă atinge pragul dublu
(cantitate $\geq$ 5 și rată $\geq$ 30\%) și emite comanda de pivotare.

\subsection{Scenariul 2: stabilirea unui canal C2 cu Sliver}

Sliver este un framework C2 open-source modern, similar conceptual cu
Cobalt Strike, folosit în testele de penetrare și în campanii reale de
post-exploatare. Scenariul a constat din descărcarea implantului Sliver
pe gazda compromisă, marcarea executabilă și pornirea acestuia; implantul
a deschis o conexiune outbound persistentă către serverul C2, iar
operatorul a emis remote comenzi standard de recon.

\begin{table}[h]
  \centering
  \caption{Scenariul 2: indicatori observați la setup Sliver}
  \label{tab:cap6-scen2}
  \renewcommand{\arraystretch}{1.08}
  \begin{tabular}{|p{0.42\textwidth}|p{0.48\textwidth}|}
    \hline
    \textbf{Acțiunea atacatorului} & \textbf{Indicator KernelTrap} \\
    \hline\hline
    Fetch implant (\texttt{curl}/\texttt{wget})              & \texttt{is\_network\_tool=1}, \texttt{has\_url\_or\_ip=1} \\
    \hline
    Pornire implant cu nume custom                          & \texttt{processName\_freq=0} (necunoscut datelor BETH) \\
    \hline
    Conexiune persistentă către C2 (\texttt{connect}, \texttt{sendto}, \texttt{recvfrom})  & \texttt{syscall\_rate\_5s} susținut pe socket non-standard \\
    \hline
    Comenzi remote: \texttt{id}, \texttt{sudo -l}, \texttt{whoami}  & \texttt{is\_basic\_recon=1}, \texttt{is\_privileged\_helper=1} \\
    \hline
  \end{tabular}
\end{table}

Rezultat: combinația dintre processName cu frecvență zero, indicatorii de
rețea și comenzile remote scor-uite ca recon împinge fereastra peste
ambele praguri. Pivotarea automată redirecționează sesiunea SSH activă
către containerul honeypot, iar regula \texttt{iptables} adăugată pe IP-ul
C2 blochează stabilirea oricăror conexiuni noi de la același atacator.

\subsection{Scenariul 3: acces la fișiere canary post-pivotare}

După redirecționarea în honeypot, atacatorul (sau implantul automat al
acestuia) explorează directorul home al utilizatorului oglindit și
încearcă să citească fișiere ce par să conțină credențiale: chei AWS
(\texttt{.aws/credentials}), cheie privată SSH (\texttt{.ssh/id\_rsa}),
fișier de mediu cu credențiale de bază de date (\texttt{.env}), note interne
(\texttt{notes.txt}) și un token de CI/CD (\texttt{.jenkins\_token}).
Pentru fiecare astfel de acces, regulile \texttt{auditd} instalate de
\texttt{provision-user.sh} emit câte un eveniment, captat de
\texttt{hp-monitor.py} și publicat pe stream-ul Redis al containerului.

\begin{table}[h]
  \centering
  \caption{Scenariul 3: indicatori observați la accesul canary}
  \label{tab:cap6-scen3}
  \renewcommand{\arraystretch}{1.08}
  \begin{tabular}{|l|l|}
    \hline
    \textbf{Fișier canary accesat}     & \textbf{Indicator KernelTrap} \\
    \hline\hline
    \texttt{\textasciitilde/.aws/credentials} & \texttt{has\_canary\_path=1}, regulă \texttt{audit\_001} \\
    \hline
    \texttt{\textasciitilde/.ssh/id\_rsa}     & \texttt{has\_canary\_path=1}, regulă \texttt{audit\_002} \\
    \hline
    \texttt{\textasciitilde/.env}              & \texttt{has\_canary\_path=1}, regulă \texttt{audit\_003} \\
    \hline
    \texttt{\textasciitilde/notes.txt}         & regulă \texttt{audit\_004}  \\
    \hline
  \end{tabular}
\end{table}

Rezultat: spre deosebire de detecția probabilistică din scenariile
precedente, accesul la fișiere canary constituie un indicator binar de
compromitere, niciun utilizator legitim nu are motiv să deschidă aceste
fișiere, întrucât conțin credențiale fictive plasate intenționat ca
momeală. Indicatorii apar imediat în panoul de control și confirmă
post-factum că pivotarea a fost justificată, oferind în plus telemetrie
asupra credențialelor pe care atacatorul le urmărește.

% ----------------------------------------------------------------
\section{Limitări}
% ----------------------------------------------------------------

Această secțiune enunță explicit clasele de atac pe care sistemul nu le
acoperă eficient. Scopul este de a delimita responsabilitatea KernelTrap și
de a evita așteptări nerealiste asupra protecției oferite.

\textbf{Atacatori avansați cu unelte proprii.} Un actor de tip APT
(\textit{Advanced Persistent Threat}, atacator cu resurse și obiective de
nivel statal sau corporativ) care folosește binare scrise sau modificate
de el însuși poate ocoli o parte din indicatorii sistemului. Două categorii
de tehnici sunt în mod special problematice. Prima: nume de proces alese
deliberat pentru a semăna cu daemoni legitimi de sistem (de exemplu, un
implant numit \texttt{kworker/0:1} sau \texttt{systemd-resolved-helper})
păcălește codificarea prin frecvență a câmpului \texttt{processName} și
lexicoanele de procese cunoscute. A doua: tehnici care nu lasă deloc urma
caracteristică a unei comenzi noi, adică cod executat direct din memorie, fără
ca un proces nou să fie creat prin syscall-ul \texttt{execve} (\textit{fileless});
injectarea unei biblioteci dinamice ostile într-un proces deja existent și
benign (\textit{DLL/SO injection} în \texttt{sshd} sau \texttt{bash});
sau interpunerea pe syscall-urile unui proces de încredere prin
\texttt{ptrace} (un mecanism de depanare folosit pentru a controla execuția
altui proces). Aceste tehnici sunt caracteristice operatorilor cu resurse
considerabile și nu constituie profilul de atacator țintă al sistemului.

\textbf{Atacuri pre-autentificare.} Agentul monitorizează exclusiv UID-urile
cu sesiuni SSH interactive active (terminalele virtuale \texttt{pts/*}).
În consecință, ies complet din raza sa de acoperire: brute-force-ul SSH
(încercări automate de parole care nu produc o sesiune validă), exploit-urile
care vizează direct procesul \texttt{sshd} înainte de stabilirea sesiunii,
precum și atacurile asupra altor servicii expuse public (servere web,
baze de date).

\textbf{Exploit de kernel sau orbire a Tracee.} Un atac care reușește să
dezactiveze sau să modifice programele eBPF deja încărcate în kernel orbește
complet agentul. Aceasta poate veni dintr-o vulnerabilitate
\textit{zero-day} în verificatorul eBPF (componenta de kernel care validează
programele eBPF înainte de încărcare), din exploatarea capabilităților
privilegiate \texttt{CAP\_BPF} sau \texttt{CAP\_SYS\_ADMIN}, sau din
înlocuirea pur și simplu a binarului Tracee. Toate aceste tehnici necesită
oricum drepturi apropiate de \texttt{root}, situație în care protecția
oferită se reduce la asumarea explicită din secțiunea~\ref{tab:cap6-trust}
(contul \texttt{root} este considerat necompromis).

\textbf{Vectori de atac în afara SSH.} Execuția de cod la distanță printr-o
aplicație web compromisă (\textit{Remote Code Execution} pe Apache, Tomcat,
servere de aplicații), sarcini malițioase rulate prin \texttt{cron}, sau
exploit-uri ale unor servicii care rulează sub conturi de sistem alb-listate
(\texttt{www-data} pentru servere web, \texttt{mysql}, \texttt{postgres}
pentru baze de date) sunt invizibile filtrelor agentului, întrucât nu trec
prin nicio sesiune SSH. KernelTrap este construit pentru scenariul „cont
SSH neprivilegiat compromis", nu pentru compromiterea generală a tuturor
serviciilor expuse.

\subsection{Domeniul de aplicabilitate optim}

Pentru a clarifica unde KernelTrap oferă valoare maximă, tabelul
\ref{tab:cap6-fit} cartografiază principalele profile de atacator
împotriva nivelului de acoperire al sistemului.

\begin{table}[h]
  \centering
  \caption{Profile de atacator și acoperirea KernelTrap}
  \label{tab:cap6-fit}
  \renewcommand{\arraystretch}{1.08}
  \begin{tabular}{|p{0.34\textwidth}|c|p{0.42\textwidth}|}
    \hline
    \textbf{Profil atacator} & \textbf{Acoperire} & \textbf{Motiv} \\
    \hline\hline
    Script kiddies / atacatori oportuniști                 & Optimă   & Folosesc kit-uri publice (LinPEAS, LinEnum, pspy, module recon Metasploit) ușor recunoscute prin lexicoane. \\
    \hline
    Botneturi care compromit SSH prin credential stuffing  & Optimă   & După obținerea foothold-ului, execută în lanț scripturi automate de post-exploatare (recon, descărcare payload, persistență, lateral movement); generează rafale dense de syscall-uri și procese cu profil zgomotos, ușor de prins de fereastra glisantă. \\
    \hline
    Mișcare laterală prin cont SSH compromis              & Optimă   & Scenariul de proiectare; întreaga arhitectură este construită pentru acest caz. \\
    \hline
    APT cu tooling custom / fileless / 0-day kernel        & Limitată & Indicatorii bazați pe nume de proces și syscall-uri standard nu rețin profilul. \\
    \hline
    Brute-force / exploit pre-autentificare                & Niciuna  & Filtrele pe UID SSH activ exclud această fază. \\
    \hline
    RCE prin aplicații web sau servicii expuse             & Niciuna  & UID-urile sistemice sunt alb-listate; vectorul nu trece prin SSH. \\
    \hline
  \end{tabular}
\end{table}

KernelTrap își propune să ridice semnificativ costul atacatorilor de nivel
mediu segmentul în care se încadrează majoritatea atacurilor observate în
trafic real, conform sondajelor de threat intelligence fără a pretinde
acoperirea spectrului complet pe care un EDR enterprise îl tratează.

% ----------------------------------------------------------------
\section{Sinteză evaluare}
% ----------------------------------------------------------------

Evaluarea experimentală a confirmat că sistemul atinge ținta de
overhead stabilită în capitolul 5 (sub 2,5\% pe CPU pe gazda monitorizată)
și detectează cu succes atât tooling-ul de recon automat (LinPEAS), cât și
stabilirea unui canal C2 modern (Sliver), declanșând în fiecare caz
pivotarea automată a sesiunii SSH compromise.
Indicatorii binari de tip canary, observați după redirecționarea în
honeypot, oferă confirmarea independentă a faptului că pivotarea a fost
justificată și telemetrie suplimentară asupra obiectivelor atacatorului.
Limitările identificate actori avansați cu tooling custom, atacuri
pre-autentificare, vectori în afara SSH sunt asumate explicit și
poziționează KernelTrap ca un sistem de detecție și răspuns activ
specializat pe atacatori autentificați de nivel mediu, complementar
soluțiilor de protecție perimetrale și EDR-urilor generale, nu un
înlocuitor pentru acestea.
